\subsubsection{从自对偶压缩到一变量 \texorpdfstring{$\Xi$}{Xi}：严格函数方程与临界线实值性}\label{subsec:xi-from-self-dual}
本节把附录中“自对偶 + 完成化 + 单位圆实值化”的三件套压缩为一个一变量 \(\Xi\)-对象，并将函数方程理解为完成化投影系统上的自同构在截面上的自然性约束；相应地，“临界线”不再被视作 \(s\)-平面上的对称轴，而被解释为参数侧的 unitary slice（\(|u|=1\)）这一协议可见切片。

\paragraph{自对偶条件}
设 \(\Delta(z,u)\) 是二变量多项式/解析函数（在本文的核口径下典型为
\(\Delta(z,u)=\det(I-zB(u))\)，见 \S\ref{subsec:zeta-online-kernel}），并满足二变量自对偶恒等式
\[
\boxed{\ \Delta(z,u)=\Delta(uz,1/u)\ }.
\]
对同步核加权情形，该恒等式在附录 \ref{app:kernel-self-dual} 给出；并由此导出完成化偶性与系数回文约束（例如推论 \ref{cor:delta-coeff-palindrome}）。

\begin{corollary}[自对偶诱导的尺度反演：谱半径与压力的反演不变量]\label{cor:xi-selfdual-pressure-inversion}
\leanverified{paper\_xi\_selfdual\_pressure\_inversion}
进一步假设 \(\Delta(z,u)=\det(I-zB(u))\) 来自某个有限维矩阵族 \(B(u)\)。
对 \(u>0\) 定义谱半径 \(\lambda(u):=\rho(B(u))\) 与压力 \(P(u):=\log\lambda(u)\)。
则二变量自对偶恒等式推出精确反演律
\[
\lambda(u)=u\,\lambda(1/u),
\qquad
P(u)=\log u+P(1/u),
\]
等价地，中心化压力 \(P(u)-\tfrac12\log u\) 在对合 \(u\leftrightarrow 1/u\) 下不变。
\end{corollary}
\begin{proof}
对固定 \(u\)，记 \(B(u)\) 的谱为 \(\mathrm{Spec}(B(u))\)（按代数重数计）。
由 \(\Delta(z,u)=\det(I-zB(u))=\prod_{\mu\in\mathrm{Spec}(B(u))}(1-z\mu)\) 与自对偶恒等式
\(\Delta(z,u)=\Delta(uz,1/u)=\prod_{\nu\in\mathrm{Spec}(B(1/u))}(1-uz\nu)\)
可知谱满足集合恒等式（按重数）
\(
\mathrm{Spec}(B(u))=u\cdot \mathrm{Spec}(B(1/u))
\)。
取模长最大值得 \(\lambda(u)=u\,\lambda(1/u)\)；取对数即得 \(P(u)=\log u+P(1/u)\)。
最后一条等价改写为
\(
P(u)-\tfrac12\log u=P(1/u)+\tfrac12\log u
\)。
\end{proof}

\paragraph{新的一变量对象：任意底 \texorpdfstring{$L>1$}{L>1}}
取任意常数底 \(L>1\)，定义
\[
z_L(s):=L^{-s},\qquad u_L(s):=L^{2s-1},
\]
并定义一变量函数
\[
\boxed{\ \Xi_{\Delta,L}(s):=\Delta\!\bigl(z_L(s),u_L(s)\bigr)=\Delta\!\bigl(L^{-s},L^{2s-1}\bigr)\ }.
\]
当取 \(L=\lambda=\rho(B(1))\) 时，这与附录 \ref{app:kernel-functional-equation} 的截面 \(\Xi(s)\) 一致（命题 \ref{prop:kernel-xi-functional-equation}）。

\begin{proposition}[严格函数方程：投影系统自同构的截面自然性]\label{prop:xi-delta-L-functional-equation}
\leanverified{paper\_xi\_delta\_l\_functional\_equation}
在上述自对偶条件下，对任意 \(L>1\) 都有严格恒等式
\[
\boxed{\ \Xi_{\Delta,L}(s)=\Xi_{\Delta,L}(1-s)\ }.
\]
\end{proposition}
\begin{proof}
令 \(\tau:(z,u)\mapsto(uz,u^{-1})\)。自对偶恒等式 \(\Delta(z,u)=\Delta(uz,1/u)\) 等价于 \(\Delta=\Delta\circ\tau\)。
对任意 \(L>1\) 直接计算得
\[
u_L(1-s)=u_L(s)^{-1},\qquad z_L(1-s)=u_L(s)\,z_L(s),
\]
从而 \((z_L(1-s),u_L(1-s))=\tau(z_L(s),u_L(s))\)。
因此
\[
\Xi_{\Delta,L}(s)
=\Delta\!\bigl(z_L(s),u_L(s)\bigr)
=\Delta\!\Bigl(\tau\bigl(z_L(s),u_L(s)\bigr)\Bigr)
=\Delta\!\bigl(z_L(1-s),u_L(1-s)\bigr)
=\Xi_{\Delta,L}(1-s).
\]
\end{proof}

\begin{remark}[投影系统自同构：函数方程的对象]\label{rem:xi-projection-system-automorphism}
自对偶恒等式 \(\Delta(z,u)=\Delta(uz,1/u)\) 可视为二变量接口上的一个 involution：
定义
\[
\tau:\ (z,u)\longmapsto (uz,u^{-1}),
\qquad \tau^2=\mathrm{id}.
\]
则自对偶条件等价于 \(\Delta=\Delta\circ\tau\)，即 \(\Delta\) 在完成化投影系统的自同构 \(\tau\) 下不变。
\par
对任意 \(L>1\)，完成化截面映射 \(s\mapsto (z_L(s),u_L(s))\) 与 \(s\mapsto 1-s\) 在该接口上严格对齐：
\[
(z_L(1-s),u_L(1-s))
=(L^{s-1},L^{1-2s})
\;=\;
\bigl(u_L(s)\,z_L(s),\,u_L(s)^{-1}\bigr)
=\tau\bigl(z_L(s),u_L(s)\bigr).
\]
因此，命题 \ref{prop:xi-delta-L-functional-equation} 的函数方程并非“关于某条几何轴的对称”，而是 \(\tau\)-不变性在完成化截面上的拉回。该口径也解释了为何 unitary-slice 核验协议自然锁定 \(|u|=1\) 的可见域：这是一条与 \(\tau\) 相容、并支撑 Hardy 型实值化与谱穿越审计的参数切片。
\end{remark}

\paragraph{完成化与“临界线实值性”}
写
\[
\Delta(z,u)=\sum_{j=0}^{N}c_j(u)\,z^j.
\]
若系数满足完成化回文结构 \(c_j(u)=u^j c_j(1/u)\)（推论 \ref{cor:delta-coeff-palindrome}），则可定义完成化系数
\[
d_j(u):=u^{-j/2}c_j(u).
\]
附录推论 \ref{cor:completion-real-line} 表明：当 \(|u|=1\) 时，\(d_j(u)\in\RR\)（完成化后“单位圆上系数为实”）。

\begin{proposition}[临界线实值性：完成化 \(\Rightarrow\) Hardy 型实值化]\label{prop:xi-delta-L-critical-line-real}
\leanverified{paper\_xi\_delta\_l\_critical\_line\_real}
在自对偶条件且满足推论 \ref{cor:completion-real-line} 的条件下，对任意 \(L>1\) 与任意实数 \(t\) 有
\[
\boxed{\ \Xi_{\Delta,L}\!\left(\tfrac12+it\right)\in\RR\ }.
\]
\end{proposition}
\begin{proof}
取 \(s=\tfrac12+it\)，则 \(u_L(s)=L^{2it}\) 满足 \(|u_L(s)|=1\)，并且
\[
\Xi_{\Delta,L}(s)
=\Delta(L^{-s},u_L(s))
=\sum_{j=0}^{N}c_j(u_L(s))\,L^{-js}
=\sum_{j=0}^{N}d_j(u_L(s))\,\bigl(L^{-1/2}\bigr)^{j}\,e^{-ijt\log L}.
\]
由 \(|u_L(s)|=1\Rightarrow d_j(u_L(s))\in\RR\)，上述表达是一个关于实变量 \(e^{-it\log L}\) 的“实系数”三角多项式在 \(t\) 处的取值，因此其值为实数。等价地，可把附录中 Hardy 实值化 \(Z(t)\) 的证明（命题 \ref{prop:kernel-hardy-real-even}）中的 \(\log\lambda\) 替换为 \(\log L\)。
\end{proof}

\begin{remark}[临界线 \texorpdfstring{$\Leftrightarrow$}{<->} 单位圆切片]\label{rem:xi-critical-line-unit-circle}
由 \(u_L(s)=L^{2s-1}\) 可见：直线 \(\Re(s)=\tfrac12\) 恰对应 \(|u_L(s)|=1\)。
因此在本文的可核验语义中，“临界线”并非额外几何公理，更不是 \(s\)-平面上的对称轴；它是完成化投影系统在参数侧的 unitary slice，并在协议层被固定为可见域 \(X_{\mathrm{vis}}=\mathsf{Phase}\) 的唯一连续切片（定义 \ref{def:xi-visible-read-null}）。
结合注 \ref{rem:xi-projection-system-automorphism}，函数方程的结构对象是投影系统自同构 \(\tau\)，而 unitary slice 则是该自同构在可见域投影下保持稳定的审计切片：在该切片上完成化系数可取实，从而支持 Hardy 型实值化与谱穿越接口。
\end{remark}

\begin{proposition}[unitary slice 的行列式重写与谱穿越判据]\label{prop:xi-unitary-slice-spectral-crossing}
\leanverified{paper\_xi\_unitary\_slice\_spectral\_crossing}
在自对偶条件下，进一步假设
\(
\Delta(z,u)=\det(I-zB(u))
\)
来自某个有限维（可审计）矩阵族 \(B(u)\)（例如同步核加权情形，见 \S\ref{subsec:zeta-online-kernel} 与附录 \ref{app:kernel-self-dual}）。
固定 \(L>1\)，令 \(\theta:=2t\log L\)。
则对任意 \(t\in\RR\) 有严格恒等式
\[
\boxed{\;
\Xi_{\Delta,L}\!\left(\tfrac12+it\right)
=
\det\!\Bigl(I-U_L(\theta)\Bigr),\qquad
U_L(\theta):=L^{-1/2}e^{-i\theta/2}\,B(e^{i\theta})\; }.
\]
因此临界线零点满足等价判据
\[
\boxed{\;
\Xi_{\Delta,L}\!\left(\tfrac12+it\right)=0
\quad\Longleftrightarrow\quad
1\in\mathrm{Spec}\bigl(U_L(\theta)\bigr),\ \ \theta=2t\log L\; }.
\]
\end{proposition}
\begin{proof}
取 \(s=\tfrac12+it\)。则
\(
u_L(s)=L^{2it}=e^{i\theta}
\)
其中 \(\theta=2t\log L\)，并且
\(
z_L(s)=L^{-s}=L^{-1/2}e^{-it\log L}=L^{-1/2}e^{-i\theta/2}
\)。
由 \(\Xi_{\Delta,L}(s)=\Delta(z_L(s),u_L(s))=\det(I-z_L(s)B(u_L(s)))\) 得到所示重写。
零点判据由行列式为零当且仅当 \(1\) 为 \(U_L(\theta)\) 的特征值立得。证毕。
\end{proof}

\begin{remark}[可审计的谱穿越口径：从多项式根到特征值轨道]\label{rem:xi-unitary-slice-operator-audit}
命题 \ref{prop:xi-unitary-slice-spectral-crossing} 把 Hardy 截面上的零点检测从
“一元多项式 \(P_L(S)\) 的根”（命题 \ref{prop:xi-rh-interval-criterion}）
等价转写为
“unitary 参数 \(\theta\) 下归一化核 \(U_L(\theta)\) 的谱是否穿过 \(1\)”。
这一口径的意义在于：它把“临界线几何”从 \(s\)-平面的隐式条件，转化为 \(|u|=1\) 上的显式算子族谱问题，从而允许直接以数值线性代数在 \(\theta\)-轴上进行独立审计。
\par
同步核加权例子中，脚本 \path{scripts/exp_sync_kernel_unitary_slice_spectral_crossing.py} 在离散网格上同时计算
\(\det(I-U_L(\theta))\) 与 \(\min_{\lambda\in\mathrm{Spec}(U_L(\theta))}|\lambda-1|\)，并输出为可复现表格：
\begin{table}[H]
\centering
\scriptsize
\setlength{\tabcolsep}{6pt}
\caption{Unitary-slice spectral-crossing audit for the weighted sync-kernel. We scan $\theta\in[0,4\pi)$ on a uniform grid and evaluate $\Xi(\tfrac12+it)=\det(I-U_L(\theta))$ with $U_L(\theta)=L^{-1/2}e^{-i\theta/2}B(e^{i\theta})$ (here $L=3$). We report the smallest $|\det(I-U)|$, the smallest spectral distance $\min_{\lambda\in\mathrm{Spec}(U)}|\lambda-1|$, and two consistency checks: the maximal imaginary leakage of $\det(I-U)$ and the maximal deviation from the completed polynomial form $\widehat\Delta(L^{-1/2},2\cos(\theta/2))$.}
\label{tab:sync_kernel_unitary_slice_spectral_crossing}
\begin{tabular}{r r r r r}
\toprule
$N_\theta$ & $\min|\det(I-U)|$ & $\min\min|\lambda-1|$ & $\max|\Im\det|$ & $\max|\det-\widehat\Delta|$\\
\midrule
4096 & 4.562e-05 & 3.162e-05 & 5.218e-16 & 8.461e-16\\
\bottomrule
\end{tabular}
\end{table}

\end{remark}

\begin{theorem}[Hardy 截面零点的谱流计数：\texorpdfstring{$\det(I-U_L(\theta))$}{det(I-U)} 的交叉数与核维数]\label{thm:xi-hardy-slice-spectral-flow-counting}
\leanverified{paper\_xi\_hardy\_slice\_spectral\_flow\_counting}
沿用命题 \ref{prop:xi-unitary-slice-spectral-crossing} 的记号，设 \(\theta\mapsto U_L(\theta)\) 为一条在迹类范数下 \(C^1\) 的幺正算子族，并设 \(\det(I-U_L(\theta))\) 在闭区间 \([\theta_0,\theta_1]\) 内的零点离散且无聚点。
则 \([\theta_0,\theta_1]\) 内的零点总数（按代数重数计）
\[
\#\{\theta\in[\theta_0,\theta_1]:\ \det(I-U_L(\theta))=0\}
\]
等于谱流 \(\mathrm{sf}\bigl(U_L(\theta);\theta_0\to\theta_1\bigr)\) 经过点 \(1\) 的总交叉数（按交叉重数计）。
并且对每个零点 \(\theta_\ast\)，其零点重数等于
\[
\dim\ker\bigl(I-U_L(\theta_\ast)\bigr).
\]
\end{theorem}
\begin{proof}
对迹类连续的幺正族，Fredholm 行列式满足：
\(\det(I-U_L(\theta))=0\) 当且仅当 \(1\in \mathrm{Spec}(U_L(\theta))\)，并且零点重数等于 \(1\) 的特征空间维数（幺正情形下代数/几何重数一致）。
在“零点离散且无聚点”的假设下，可在每个零点邻域用谱投影的连续性选取穿越 \(1\) 的有限条特征值分支；谱流的定义正是这些分支穿越 \(1\) 的交叉数之和，因此得到计数恒等式与重数识别。证毕。
\end{proof}

\paragraph{自对偶截面上的完成化常数：\texorpdfstring{$w$}{w} 固定，\texorpdfstring{$\Xi$}{Xi} 退化为余弦多项式}
为避免与复变量 \(s\) 混淆，本段把 completion 的迹坐标记为 \(S\)。
若自对偶 \(\Delta(z,u)\) 还可写成某个完成化行列式
\(\widehat\Delta(w,S)\) 的形式
\[
\Delta(z,u)=\widehat\Delta\!\left(z\sqrt{u},\ \sqrt{u}+\frac{1}{\sqrt{u}}\right)
\]
其中 \(\widehat\Delta\) 在自对偶核上由不变量子环强制存在；例如同步核加权情形见命题 \ref{prop:sync-kernel-weighted-hat-delta}。
则沿自对偶 Dirichlet 截面 \(z=L^{-s},\ u=L^{2s-1}\) 有
\[
w:=z\sqrt{u}=L^{-s}\cdot L^{s-1/2}=L^{-1/2}\quad(\text{常数}),
\qquad
S_L(s):=\sqrt{u}+\frac{1}{\sqrt{u}}=L^{s-1/2}+L^{1/2-s},
\]
从而得到一维压缩恒等式
\[
\boxed{\ \Xi_{\Delta,L}(s)=\widehat\Delta\!\bigl(L^{-1/2},\,S_L(s)\bigr)\ }.
\]
特别地在临界线 \(s=\tfrac12+it\) 上，
\[
S_L\!\left(\tfrac12+it\right)=L^{it}+L^{-it}=2\cos(t\log L),
\]
因此 Hardy 型截面 \(\Xi_{\Delta,L}(\tfrac12+it)\) 是一个“余弦的实多项式”。

\paragraph{Lee--Yang 折叠与素相位提升}
本段把 unitary slice 上的余弦变量 \(S=2\cos(\theta/2)\) 进一步压缩到 Lee--Yang 奇异环坐标 \(y=J(e^{\mathrm{i}\theta})=-1/S^2\)；
并将 \(\theta=2t\log L\) 的相位驱动以素数分解的方式提升到算术奇异环 \(\mathfrak{R}_{S(L)}\) 的一参数流上。

\begin{proposition}[Lee--Yang 折叠与无符号化根等价]\label{prop:xi-leyang-folding-signed-roots}
令 \(F(x)\in\CC[x]\) 为多项式，并定义其 Lee--Yang 折叠
\[
F^{\sharp}(w):=F(\sqrt{w})\,F(-\sqrt{w})\in\CC[w],
\]
其中乘积保证 \(F^{\sharp}\) 为 \(w\) 的多项式且不依赖平方根分支。
则对任意实数 \(x\) 有
\[
F(x)=0\quad\Longleftrightarrow\quad F^{\sharp}(x^2)=0,
\]
并且 \(F^{\sharp}(x^2)=0\) 的根分支与 \(F\) 在 \(\pm x\) 处的零点对应。
\leanverified{paper\_xi\_leyang\_folding\_signed\_roots}
\end{proposition}
\begin{proof}
恒等式 \(F^{\sharp}(x^2)=F(x)F(-x)\) 直接给出零点等价与分支对应。证毕。
\end{proof}

\begin{proposition}[单位圆扫描到奇异环扫描的代数折叠]\label{prop:xi-unitary-slice-leyang-folded-scan}
固定 \(L>1\)，并令 \(w:=L^{-1/2}\)。
沿用本节的完成化表示，置
\[
P_L(S):=\widehat\Delta(w,S)\in\RR[S],
\qquad
P_L^{\sharp}(X):=P_L(\sqrt{X})P_L(-\sqrt{X})\in\RR[X].
\]
令 \(\theta\in\RR\) 并置 \(z=e^{\mathrm{i}\theta/2}\in\TT\)。
则 Lee--Yang 坐标
\[
y:=J(z^2)=-\frac{1}{(z+z^{-1})^2}=-\frac{1}{(2\cos(\theta/2))^2}\in(-\infty,-\tfrac14]
\]
满足 \(X:=(2\cos(\theta/2))^2=-1/y\)，并且在 unitary slice 上的零点事件满足等价判别
\[
P_L\!\bigl(2\cos(\theta/2)\bigr)=0
\quad\Longleftrightarrow\quad
P_L^{\sharp}(-1/y)=0.
\]
因此，\(\theta\)-轴上的单位圆扫描可被严格转写为 Lee--Yang 实射线坐标 \(y\in(-\infty,-\tfrac14]\) 上的实根审计问题。
\end{proposition}
\begin{proof}
由定义 \(X=(2\cos(\theta/2))^2\) 且 \(y=-1/X\)。
命题 \ref{prop:xi-leyang-folding-signed-roots} 对 \(F=P_L\) 给出
\(
P_L(2\cos(\theta/2))=0\iff P_L^{\sharp}(X)=0
\)；
再以 \(X=-1/y\) 代换即得结论。证毕。
\end{proof}

\begin{theorem}[幺正切片的素相位提升与算术奇异环流]\label{thm:xi-unitary-slice-prime-phase-lift-to-arithmetic-singular-ring}
令 \(L\ge 2\) 为整数，并记其素因子集合 \(S(L)\)，写
\[
L=\prod_{p\in S(L)}p^{v_p(L)}.
\]
对 \(t\in\RR\) 定义全局相位与局部素相位
\[
u_L(t):=L^{2\mathrm{i}t}=\exp\!\bigl(2\mathrm{i}t\log L\bigr)\in\TT,
\qquad
u_{p,L}(t):=\exp\!\bigl(2\mathrm{i}t\,v_p(L)\log p\bigr)\in\TT\quad(p\in S(L)).
\]
则有严格分解
\[
u_L(t)=\prod_{p\in S(L)}u_{p,L}(t),
\qquad
\sum_{p\in S(L)}v_p(L)\log p=\log L.
\]
从而存在一条连续一参数子群
\[
\tau^{\mathrm{prime}}_{L}:\RR\longrightarrow \mathfrak{R}_{S(L)}=\Sigma_{S(L)}\times_{\TT}\TT^{S(L)}
\]
使得其 \(\TT^{S(L)}\) 分量为 \(\bigl(u_{p,L}(t)\bigr)_{p\in S(L)}\)，并且其 \(\Sigma_{S(L)}\) 分量在 \(\pi_{S(L)}\) 下投影为 \(u_L(t)\)。
\end{theorem}
\begin{proof}
由对数分解恒等式立得 \(\prod_{p\in S(L)}u_{p,L}(t)=u_L(t)\)，故曲线
\(
t\mapsto (x(t),(u_{p,L}(t))_{p\in S(L)})
\)
只需取 \(x(t)\in\Sigma_{S(L)}\) 满足 \(\pi_{S(L)}(x(t))=u_L(t)\) 即可落入纤维积 \(\mathfrak{R}_{S(L)}\)。
另一方面，定理 \ref{thm:cdim-arithmetic-singular-ring-one-parameter-subgroups} 给出一参数子群与频率向量的自然对应：
取 \(\omega_{p,L}:=v_p(L)\log p\)，则 \(\sum_{p\in S(L)}\omega_{p,L}=\log L\) 且 \(\mathrm{pr}_2(\tau^{\mathrm{prime}}_L(t))=(e^{\mathrm{i}\omega_{p,L}t})_{p\in S(L)}=(u_{p,L}(t))_{p\in S(L)}\)。
由该对应的唯一性，一参数子群 \(\tau^{\mathrm{prime}}_L\) 亦唯一确定，并满足所述投影条件。证毕。
\end{proof}

\begin{conjecture}[算术奇异环表示的幺正切片因子化与正性约束]\label{conj:xi-unitary-slice-arithmetic-singular-ring-representation}
对每个整数尺度 \(L\ge 2\)，存在一个与 \(\theta\mapsto U_L(\theta)\)（命题 \ref{prop:xi-unitary-slice-spectral-crossing}）兼容的连续酉表示
\[
\rho_L:\mathfrak{R}_{S(L)}\to \mathrm{U}(D_L)
\]
使得对所有 \(t\in\RR\) 成立
\[
U_L(2t\log L)=\rho_L\bigl(\tau^{\mathrm{prime}}_L(t)\bigr),
\]
其中 \(\tau^{\mathrm{prime}}_L\) 为定理 \ref{thm:xi-unitary-slice-prime-phase-lift-to-arithmetic-singular-ring} 所构造的素相位流。
进一步，沿 Lee--Yang 坐标 \(y=J(e^{\mathrm{i}\theta})\in(-\infty,-\tfrac14]\) 的折叠审计（命题 \ref{prop:xi-unitary-slice-leyang-folded-scan}）下，
由行列式观测 \(y\mapsto \det(I-\rho_L(\cdot))\) 所诱导的实轴函数属于 Herglotz--Nevanlinna（或等价的 Stieltjes）正性类，从而与本文的 Carath\'eodory/Toeplitz 证书接口及 Stieltjes--Hankel 分层正性口径一致（见 \S\ref{subsubsec:xi-axis-swap-stieltjes-hankel-positivity}）。
\end{conjecture}

\endinput


在完成化压缩口径下，越界根对应的“径向偏离量”可被写成一个完全显式的 \(\mathrm{arcosh}\) 闭式；该闭式把离线竖线的水平偏移从叙事性描述压缩为由 \(P_L\) 根位置刚性决定的数值不变量。\footnote{据内部文稿 \cite{Internal2025ResolutionFoldingPhiPi}（2025）在“相位坐标门—有限窗审计—orbit--Levi 刚性”并置且不增设额外公理的口径下，该偏离量与常数层规范归一化一起形成可独立审计的封闭链。}

\begin{corollary}[单尺度完成化的径向偏离量闭式：越界根的离线距离由 \texorpdfstring{$\mathrm{arcosh}$}{arcosh} 刚性给出]\label{cor:xi-single-scale-radial-deviation-arcosh}
\leanverified{paper\_xi\_single\_scale\_radial\_deviation\_arcosh}
在命题 \ref{prop:xi-rh-interval-criterion} 的完成化表示下，令
\[
\Xi_{\Delta,L}(s)=P_L\!\bigl(S_L(s)\bigr),\qquad S_L(s)=L^{s-\frac12}+L^{\frac12-s},
\]
其中 \(P_L\in\RR[S]\)。
取 \(y:=L^{s-\frac12}\)，则 \(S_L(s)=y+y^{-1}\)。
令 \(S_\star\in\RR\) 为 \(P_L\) 的一根（按重数计），并令 \(y_\star^{\pm}\) 为二次方程 \(y+y^{-1}=S_\star\) 的两解：
\[
y_\star^{\pm}=\frac{S_\star\pm\sqrt{S_\star^2-4}}{2}.
\]
则由复对数多值性，对每个符号 \(\pm\) 与 \(k\in\ZZ\) 有一条零点塔
\[
s=\frac12+\frac{\log y_\star^{\pm}+2\pi i k}{\log L},
\]
并且：
\begin{enumerate}
  \item 若 \(S_\star\in[-2,2]\)，则 \(|y_\star^\pm|=1\)，从而该根所诱导的全部零点满足 \(\Re(s)=\tfrac12\)；
  \item 若 \(|S_\star|>2\)，则该根所诱导的离线零点满足
  \[
  \bigl|\Re(s)-\tfrac12\bigr|
  =\frac{\mathrm{arcosh}\bigl(|S_\star|/2\bigr)}{\log L}.
  \]
\end{enumerate}
\end{corollary}

\begin{proof}
由 \(S_L(s)=y+y^{-1}\) 得二次方程 \(y^2-S_\star y+1=0\)，从而 \(y=y_\star^\pm\)。
再由 \(y=L^{s-\frac12}\) 得
\(
\log y=(s-\frac12)\log L-2\pi i k
\)
（\(k\in\ZZ\)），即得到零点塔表达。
\par
当 \(S_\star\in[-2,2]\) 时可写 \(S_\star=2\cos\theta\)，则 \(y_\star^\pm=e^{\pm i\theta}\)，故 \(|y_\star^\pm|=1\)。
当 \(|S_\star|>2\) 时令 \(|S_\star|=2\cosh\eta\)（\(\eta>0\)，即 \(\eta=\mathrm{arcosh}(|S_\star|/2)\)），则两根满足 \(|y_\star^\pm|=e^{\pm\eta}\)。
由 \(|y|=L^{\Re(s)-1/2}\)（命题 \ref{prop:xi-rh-interval-criterion} 的证明）得
\(
|\Re(s)-\tfrac12|
=\eta/\log L
\)，即结论。
\end{proof}


\begin{proposition}[零点临界线版的纯代数判据：\texorpdfstring{$\Xi$}{Xi}-RH \(\Longleftrightarrow\) 根区间]\label{prop:xi-rh-interval-criterion}
\leanverified{paper\_xi\_rh\_interval\_criterion}
在上述完成化表示下，令
\[
P_L(S):=\widehat\Delta\!\bigl(L^{-1/2},S\bigr)\in\RR[S].
\]
则下述两条陈述等价：
\begin{enumerate}
  \item \(\Xi_{\Delta,L}\) 的全部零点都满足 \(\Re(s)=\tfrac12\)；
  \item \(P_L(S)\) 的全部根都落在实区间 \([-2,2]\) 内。
\end{enumerate}
\end{proposition}
\begin{proof}
令 \(y:=L^{s-1/2}\in\CC^\times\)。则 \(S_L(s)=y+y^{-1}\)，从而
\[
\Xi_{\Delta,L}(s)=\widehat\Delta\!\bigl(L^{-1/2},\,y+y^{-1}\bigr)=P_L(y+y^{-1}).
\]
因此 \(\Xi_{\Delta,L}(s)=0\) 当且仅当 \(S_0:=y+y^{-1}\) 为 \(P_L\) 的某个根。
\par
另一方面，
\[
|y|=|L^{s-1/2}|=L^{\Re(s)-1/2},
\]
故 \(\Re(s)=\tfrac12\) 当且仅当 \(|y|=1\)。
而 \(|y|=1\) 当且仅当 \(y=e^{i\theta}\)（某个 \(\theta\in\RR\)），此时
\(
y+y^{-1}=2\cos\theta\in[-2,2]
\) 为实数。
反之，若 \(S_0\in[-2,2]\) 为实数，则方程 \(y+y^{-1}=S_0\) 的解满足 \(|y|=1\)。
因此 “所有零点都满足 \(\Re(s)=\tfrac12\)” 等价于 “\(P_L\) 的所有根都在 \([-2,2]\)”。
\end{proof}

\begin{corollary}[单位圆根判据：\texorpdfstring{$\Xi$}{Xi}-RH \(\Longleftrightarrow\) 自倒数多项式的单位圆根分布]\label{cor:xi-rh-unit-circle-reciprocal}
在命题 \ref{prop:xi-rh-interval-criterion} 的记号下，令 \(d:=\deg P_L\) 并定义
\[
\Phi_L(y):=y^{d}P_L(y+y^{-1})\in\RR[y].
\]
则：
\begin{enumerate}
  \item \(\Phi_L\) 为自倒数（reciprocal/palindromic）多项式，即
  \[
  \boxed{\ y^{2d}\Phi_L(1/y)=\Phi_L(y)\ }.
  \]
  \item 下列陈述等价：
  \begin{enumerate}
    \item[(i)] \(\Xi_{\Delta,L}\) 的全部零点都满足 \(\Re(s)=\tfrac12\)；
    \item[(ii)] \(\Phi_L(y)\) 的全部复根都满足 \(|y|=1\)；
    \item[(iii)] \(P_L(S)\) 的全部根都落在实区间 \([-2,2]\) 内。
  \end{enumerate}
\end{enumerate}
\end{corollary}
\leanverified{paper\_xi\_rh\_unit\_circle\_reciprocal}
\begin{proof}
由定义 \(\Phi_L(y)=y^{d}P_L(y+y^{-1})\)，把 \(y\mapsto 1/y\) 代入得
\(
\Phi_L(1/y)=y^{-d}P_L(y+y^{-1})
\)，
故 \(y^{2d}\Phi_L(1/y)=y^{d}P_L(y+y^{-1})=\Phi_L(y)\)，从而 \(\Phi_L\) 自倒数。
\par
另一方面，\(\Phi_L(y)=0\) 当且仅当 \(P_L(S_0)=0\) 且 \(S_0=y+y^{-1}\)。
若 \(S_0\in[-2,2]\)，则 \(y\) 必可写为 \(y=e^{i\theta}\)（\(\theta\in\RR\)），从而 \(|y|=1\)；反之若 \(|y|=1\)，则 \(S_0=y+y^{-1}\in[-2,2]\)。
因此 “\(P_L\) 的根全在 \([-2,2]\)” 等价于 “\(\Phi_L\) 的根全在单位圆”，再与命题 \ref{prop:xi-rh-interval-criterion} 合并即得 (i)\(\Leftrightarrow\)(ii)\(\Leftrightarrow\)(iii)。
\end{proof}

\begin{proposition}[Cayley 超双曲化：单位圆根判据的实轴等价]\label{prop:xi-rh-cayley-hyperbolicity}
\leanverified{paper\_xi\_rh\_cayley\_hyperbolicity}
在推论 \ref{cor:xi-rh-unit-circle-reciprocal} 的记号下，令 \(D:=\deg\Phi_L=2d\)。
定义 Cayley 变换
\begin{equation}\label{eq:xi-cayley-y-from-x}
\boxed{\;
y=\mathcal{C}(x):=\frac{x-\mathrm{i}}{x+\mathrm{i}},
\qquad
x=\mathcal{C}^{-1}(y)=\mathrm{i}\,\frac{1+y}{1-y}\;}
\end{equation}
（它把 \(\overline{\RR}=\RR\cup\{\infty\}\) 双射到单位圆 \(\TT\)，并满足 \(\mathcal{C}(\infty)=1\)）。
并定义一元多项式
\begin{equation}\label{eq:xi-cayley-psi}
\boxed{\;
\Psi_L(x):=(x+\mathrm{i})^{D}\,\Phi_L\!\left(\frac{x-\mathrm{i}}{x+\mathrm{i}}\right)\; }.
\end{equation}
则下列条件等价：
\begin{enumerate}
  \item \(\Phi_L(y)\) 的全部复根都满足 \(|y|=1\)；
  \item \(\Psi_L(x)\) 的全部复根都落在 \(\RR\cup\{\infty\}\)（按射影意义计入 \(x=\infty\) 的根）。
\end{enumerate}
当 \(\Phi_L\in\RR[y]\) 为自倒数多项式时，\(\Psi_L\) 的系数经归一化可取为实数，因此上述条件可表述为 \(\Psi_L\) 的一元实\textbf{超双曲性}（全实根）。
\end{proposition}
\begin{proof}
由 \eqref{eq:xi-cayley-y-from-x} 可直接验证：对 \(x\in\RR\)，有 \(|\mathcal{C}(x)|=1\)；反之对 \(|y|=1\) 且 \(y\neq 1\)，有 \(x=\mathrm{i}(1+y)/(1-y)\in\RR\)，并且 \(y=1\) 对应 \(x=\infty\)。
因此 Cayley 变换在根集合层面给出双射
\(
\mathrm{Roots}(\Phi_L)\subseteq\TT
\Longleftrightarrow
\mathrm{Roots}(\Psi_L)\subseteq\overline{\RR}
\)：
若 \(\Phi_L(y_0)=0\) 且 \(y_0=\mathcal{C}(x_0)\)（\(x_0\in\RR\) 或 \(x_0=\infty\)），则清分母得 \(\Psi_L(x_0)=0\)（对 \(x_0=\infty\) 以射影根计入）；反之若 \(\Psi_L(x_0)=0\) 则 \(y_0=\mathcal{C}(x_0)\) 为 \(\Phi_L\) 的根且 \(|y_0|=1\)。
\par
若 \(\Phi_L\) 为自倒数且实系数，写 \(\Phi_L(y)=\sum_{k=0}^{D}a_k y^k\) 且 \(a_k=a_{D-k}\in\RR\)。
由 \eqref{eq:xi-cayley-psi} 展开得
\[
\Psi_L(x)=\sum_{k=0}^{D} a_k\,(x+\mathrm{i})^{D-k}(x-\mathrm{i})^{k}.
\]
利用对称性 \(a_k=a_{D-k}\) 成对合并 \(k\) 与 \(D-k\) 项即可把右端写成 \(\sum_k 2a_k\,\Re\bigl((x+\mathrm{i})^{D-k}(x-\mathrm{i})^{k}\bigr)\)，从而 \(\Psi_L\)（经整体常数归一化）具有实系数。证毕。
\end{proof}

\begin{proposition}[Cohn 型导数盘内性判别：单位圆根的等价证书]\label{prop:xi-rh-cohn-derivative}
\leanverified{paper\_xi\_rh\_cohn\_derivative}
在推论 \ref{cor:xi-rh-unit-circle-reciprocal} 的记号下，\(\Phi_L\in\RR[y]\) 为自倒数多项式且 \(\Phi_L(0)=\mathrm{lc}(P_L)\neq 0\)。
则下列陈述等价：
\begin{enumerate}
  \item \(\Phi_L\) 的全部复根满足 \(|y|=1\)；
  \item \(\Phi_L\) 为 self-inversive，并且其导数 \(\Phi_L'(y)\) 的全部复根满足 \(|y|\le 1\)。
\end{enumerate}
因此 \(\Xi\)-RH 条件（命题 \ref{prop:xi-rh-interval-criterion} / 推论 \ref{cor:xi-rh-unit-circle-reciprocal} 的任一等价形式）可在 \(|y|\le 1\) 的\textbf{导数盘内性}上给出。
\end{proposition}
\begin{proof}
自倒数与实系数蕴含
\(
\Phi_L(y)=y^{D}\overline{\Phi_L(1/\overline{y})}
\)
（其中 \(D=\deg\Phi_L\)），即 \(\Phi_L\) 为 self-inversive。
由 Cohn 定理：多项式的全部零点在单位圆上当且仅当其为 self-inversive 且其导数的全部零点落在闭单位圆盘内；参见 \cite{LalinSmyth2012UnimodularitySelfInversive}（其中回顾并推广了 Cohn 的原始判别）。
\end{proof}

\begin{remark}[Schur--Cohn 半正定证书接口：把“单位圆根/区间根”转写为可审计 LMI]\label{rem:xi-schur-cohn-sdp-interface}
令 \(p(y)=a_0+a_1y+\cdots+a_n y^n\in\CC[y]\)（\(a_0\neq 0\)）。
定义下三角 Toeplitz 矩阵
\[
X(p):=\begin{pmatrix}
a_0 & 0 & \cdots & 0\\
a_1 & a_0 & \ddots & \vdots\\
\vdots & \ddots & \ddots & 0\\
a_{n-1} & \cdots & a_1 & a_0
\end{pmatrix},
\qquad
Y(p):=\begin{pmatrix}
\overline{a_n} & 0 & \cdots & 0\\
\overline{a_{n-1}} & \overline{a_n} & \ddots & \vdots\\
\vdots & \ddots & \ddots & 0\\
\overline{a_1} & \cdots & \overline{a_{n-1}} & \overline{a_n}
\end{pmatrix},
\]
并令 Schur--Cohn Hermitian 形式矩阵
\[
H(p):=X(p)X(p)^\ast-Y(p)Y(p)^\ast\in\Mat_n(\CC)_{\mathrm{Herm}}.
\]
Schur--Cohn 理论给出：若 \(p\) 在单位圆上无零点，则 \(p\) 的全部零点落在开单位圆盘 \(|y|<1\) 当且仅当 \(H(p)\succ 0\)；并且 \(H(p)\) 的惯性指数给出单位圆内/外零点计数（标量情形的综述见 \cite{DymYoung1990SchurCohnMatrixPolynomials}）。
\par
结合命题 \ref{prop:xi-rh-cohn-derivative}，在典型非退化情形（例如 \(\Phi_L\) 的单位圆根均为单根）下，\(\Xi\)-RH 等价于 \(H(\Phi_L')\succ 0\)；
而在允许重根的边界情形，可将 \(|y|\le 1\) 的闭盘条件替换为严格版 \(|y|<1\)，或将 \(H(\Phi_L')\succ 0\) 放松为半正定 \(H(\Phi_L')\succeq 0\) 并配合边界根的有限审计。
\par
当 \(\Phi_L\)（或更一般的通道多项式 \(\Phi_{L,\rho}\)）的系数落在不变量子环 \(\QQ[u+u^{-1}]\) 中时，沿单位圆 \(u=e^{it}\) 可把系数视为 \(x=2\cos t\in[-2,2]\) 的多项式；于是 \(H(\Phi_{L,\rho}')\) 成为一个\textbf{一元多项式矩阵} \(H(x)\)，从而“全角度”核版 (G)RH 可压缩为区间半正定约束
\[
H(x)\succeq 0\qquad(x\in[-2,2]).
\]
该区间半正定性进一步可通过 SOS/半定规划给出可审计证书。
\end{remark}

\begin{theorem}[视界 Carath\'eodory--Toeplitz 证书：\texorpdfstring{$\Xi$}{Xi}-切片区间根判据的统一正性框架]\label{thm:xi-slice-horizon-carath-toeplitz}
\leanverified{paper\_xi\_slice\_horizon\_carath\_toeplitz}
沿用命题 \ref{prop:xi-rh-interval-criterion} 与推论 \ref{cor:xi-rh-unit-circle-reciprocal} 的记号：
给定 \(L>1\)，令
\(
P_L(S)=\widehat\Delta(L^{-1/2},S)\in\RR[S]
\),
令 \(d=\deg P_L\)，并令
\(
\Phi_L(y)=y^{d}P_L(y+y^{-1})\in\RR[y]
\)
（于是 \(D:=\deg\Phi_L=2d\) 且 \(\Phi_L(0)\neq 0\)）。
定义其\textbf{视界 Carath\'eodory 函数}（圆盘变量记为 \(\omega\in\mathbb{D}\)）
\begin{equation}\label{eq:xi-slice-carath-def}
\boxed{\;
\mathscr{C}_{\Phi_L}(\omega)
:=
1-\frac{2}{D}\,\omega\,\frac{\Phi_L'(\omega)}{\Phi_L(\omega)},
\qquad \omega\in\mathbb{D}.
\;}
\end{equation}
将 \(\mathscr{C}_{\Phi_L}\) 在 \(\omega=0\) 处作归一化幂级数展开：
\[
\mathscr{C}_{\Phi_L}(\omega)=c_0+2\sum_{n\ge 1}c_n\omega^n,
\qquad c_{-n}:=\overline{c_n},
\]
并对 \(N\ge 0\) 定义 Toeplitz 主子矩阵
\[
\mathbf{T}_N:=\bigl(c_{j-k}\bigr)_{0\le j,k\le N}.
\]
则下列命题等价：
\begin{enumerate}
  \item \(\Xi_{\Delta,L}\) 的全部零点满足 \(\Re(s)=\tfrac12\)（等价地：\(P_L\) 的全部根落在 \([-2,2]\) 内）；
  \item \(\Phi_L\) 的全部复根满足 \(|y|=1\)；
  \item \(\mathscr{C}_{\Phi_L}\) 在 \(\mathbb{D}\) 上解析且 \(\Re \mathscr{C}_{\Phi_L}(\omega)\ge 0\)（即 \(\mathscr{C}_{\Phi_L}\) 属于 Carath\'eodory 类）；
  \item 对所有 \(N\ge 0\)，有 Toeplitz--PSD 层级 \(\mathbf{T}_N\succeq 0\)。
\end{enumerate}
\end{theorem}
\begin{proof}
\((1)\Longleftrightarrow(2)\) 已由推论 \ref{cor:xi-rh-unit-circle-reciprocal} 给出。
\par
\((2)\Rightarrow(3)\) 设 \(\Phi_L\) 的根（按重数计）为 \(\{y_j\}_{j=1}^{D}\subset\TT\)，并令 \(\alpha_j:=1/y_j=\overline{y_j}\in\TT\)。
则 \(\Phi_L(\omega)=\Phi_L(0)\prod_{j=1}^{D}(1-\alpha_j\omega)\)。
由 \eqref{eq:xi-slice-carath-def} 直接计算得
\[
\mathscr{C}_{\Phi_L}(\omega)
=1+\frac{2}{D}\sum_{j=1}^{D}\frac{\alpha_j\omega}{1-\alpha_j\omega}
=\frac{1}{D}\sum_{j=1}^{D}\frac{1+\alpha_j\omega}{1-\alpha_j\omega}.
\]
对每个 \(\alpha\in\TT\)，函数 \(\omega\mapsto\frac{1+\alpha\omega}{1-\alpha\omega}\) 在 \(\mathbb{D}\) 上解析且具有正实部；
故其平均亦在 \(\mathbb{D}\) 上解析且满足 \(\Re\mathscr{C}_{\Phi_L}\ge 0\)，即 \(\mathscr{C}_{\Phi_L}\) 属于 Carath\'eodory 类。
\par
\((3)\Rightarrow(2)\) 若 \(\mathscr{C}_{\Phi_L}\) 在 \(\mathbb{D}\) 上解析，则 \(\Phi_L\) 在 \(\mathbb{D}\) 内无零点（否则 \(\Phi_L'/\Phi_L\) 在盘内产生极点）。
而 \(\Phi_L\) 为自倒数多项式（推论 \ref{cor:xi-rh-unit-circle-reciprocal}），故其零点按 \(y\leftrightarrow 1/\overline y\) 成对对应：若存在 \(|y_0|>1\) 的零点，则 \(1/\overline y_0\in\mathbb{D}\) 亦为零点，矛盾。
因此 \(\Phi_L\) 亦无 \(|y|>1\) 的零点，从而所有零点都落在 \(\TT\) 上。
\par
\((3)\Longleftrightarrow(4)\) 沿用附录 \ref{subsubsec:app-horizon-weyl-carath-toeplitz} 的 Carath\'eodory/Herglotz 表示与截断矩完备性（Carath\'eodory--Fej\'er），
\(\Re f\ge 0\)（\(f\) 在 \(\mathbb{D}\) 上解析）当且仅当其幂级数矩生成的 Toeplitz 主子矩阵族半正定；
将该一般结论应用于 \(f=\mathscr{C}_{\Phi_L}\) 即得。证毕。
\end{proof}

\begin{corollary}[视界 Weyl/Herglotz 版本：\texorpdfstring{\(\Xi\)}{Xi}-切片的上半平面正性判据]\label{cor:xi-slice-weyl-herglotz}
沿用定理 \ref{thm:xi-slice-horizon-carath-toeplitz} 的记号，并定义完成化临界截面
\[
E_{\Delta,L}(z):=\Xi_{\Delta,L}\!\left(\tfrac12+\mathrm{i}z\right),
\qquad z\in\CC,
\]
以及其\textbf{视界 Weyl/Carath\'eodory 对数导数}
\[
m_{\Delta,L}(z):=-\frac{E'_{\Delta,L}(z)}{E_{\Delta,L}(z)},
\qquad z\in\mathbb{H}.
\]
则定理 \ref{thm:xi-slice-horizon-carath-toeplitz} 中 (1)--(4) 亦等价于：
\begin{enumerate}
  \item[(5)] \(m_{\Delta,L}\) 在 \(\mathbb{H}\) 上解析且满足 \(\Im m_{\Delta,L}(z)\ge 0\)（即 \(m_{\Delta,L}\) 为 Herglotz/Nevanlinna 函数）。
\end{enumerate}
\end{corollary}
\begin{proof}
定理 \ref{thm:xi-slice-horizon-carath-toeplitz} 的 (1) 等价于 \(\Xi_{\Delta,L}\) 的全部零点满足 \(\Re(s)=\tfrac12\)，亦即 \(E_{\Delta,L}\) 的全部零点为实数（因为 \(E_{\Delta,L}(z)=0\iff \Xi_{\Delta,L}(\tfrac12+\mathrm{i}z)=0\)）。
由命题 \ref{prop:xi-delta-L-critical-line-real} 可知 \(E_{\Delta,L}(t)\in\RR\) 对一切 \(t\in\RR\) 成立，从而 \(E_{\Delta,L}(\overline z)=\overline{E_{\Delta,L}(z)}\) 并使零点关于实轴成共轭对称。
因此，“零点全实”当且仅当 \(E_{\Delta,L}\) 在 \(\mathbb{H}\) 内无零点，等价于 \(m_{\Delta,L}\) 在 \(\mathbb{H}\) 上解析。
\par
在该解析性前提下，\(m_{\Delta,L}=-E'_{\Delta,L}/E_{\Delta,L}\) 的 Herglotz 正性与“零点全实”之间的等价性可由附录定理 \ref{thm:app-horizon-weyl-herglotz} 的同型论证（把 \(E_{\zeta}\) 替换为 \(E_{\Delta,L}\)）得到；
合并即得 (1)\(\Longleftrightarrow\)(5)，再与定理 \ref{thm:xi-slice-horizon-carath-toeplitz} 的 (1)\(\Longleftrightarrow\)(2)\(\Longleftrightarrow\)(3)\(\Longleftrightarrow\)(4) 合并即得结论。证毕。
\end{proof}

\begin{remark}[有限核的 Ramanujan--Herglotz 同型接口（完成化函数方程的作用）]\label{rem:xi-finite-kernel-ramanujan-herglotz-interface}
对有限维核的 Dirichlet 化 \(\zeta\)-函数 \(Z_M(s)=\det(I-\lambda^{-s}M)^{-1}\)，Ramanujan 半径界（kernel RH）等价于右半开带 \(\tfrac12<\Re(s)<1\) 内无非平凡极点（推论 \ref{cor:finite-rh-halfstrip-nopoles}）。
若进一步满足谱互反闭包 \(\mu\mapsto \lambda/\mu\)（命题 \ref{prop:reciprocal-spectrum-functional-equation}），则完成化函数方程把极点几何强制为关于 \(\Re(s)=\tfrac12\) 的对称结构；
于是“半条带无极点”被提升为临界带 toy-RH 的临界线定位，并可等价转写为上半平面 Weyl/Herglotz 正性判据，从而与 \(\Xi_\zeta\) 的视界 Weyl 判别（附录定理 \ref{thm:app-horizon-weyl-herglotz}）在证书语言上同型。
\end{remark}

\begin{remark}[有限维情形的视界谱测度与反射正性]\label{rem:xi-slice-horizon-measure-reflection-positivity}
在定理 \ref{thm:xi-slice-horizon-carath-toeplitz} 的等价条件成立时，
\(\mathscr{C}_{\Phi_L}\) 为有理 Carath\'eodory 函数，故其 Herglotz 表示可取为有限原子测度（见附录 \ref{subsubsec:app-horizon-weyl-carath-toeplitz} 的表示式 \eqref{eq:app-horizon-herglotz-disk} 的有理情形，或参见 \cite{Simon2005OPUC1}）：
若 \(\mathrm{Roots}(\Phi_L)=\{y_j\}_{j=1}^{D}\subset\TT\)，则存在概率测度
\(
\mu_{\Phi_L}:=\frac{1}{D}\sum_{j=1}^{D}\delta_{y_j}
\)
使
\[
\mathscr{C}_{\Phi_L}(\omega)=\int_{\TT}\frac{\xi+\omega}{\xi-\omega}\,d\mu_{\Phi_L}(\xi).
\]
其 Fourier 矩 \(c_n=\int_{\TT}\xi^{-n}\,d\mu_{\Phi_L}(\xi)\) 即为“视界多极矩/毛发系数”；
而 Toeplitz--PSD 层级 \(\mathbf{T}_N\succeq 0\) 给出该矩序列来自正测度的\emph{反射正性}证书（与附录注 \ref{rem:app-horizon-multipole-positivity-hair} 的语义对齐）。
\end{remark}

\begin{corollary}[通道化版本：核版 (G)RH 的视界 Toeplitz--PSD 证书]\label{cor:xi-channel-horizon-toeplitz-certificate}
\leanverified{paper\_xi\_channel\_horizon\_toeplitz\_certificate}
在注 \ref{rem:xi-artin-channel-pipeline} 的通道化设定下，若对每个通道 \(\rho\) 的完成化截面 \(\Xi_{\rho,L}\) 定义相应的自倒数多项式 \(\Phi_{L,\rho}\)（推论 \ref{cor:xi-rh-unit-circle-reciprocal} 的逐通道版本），并定义 \(\mathscr{C}_{\Phi_{L,\rho}}\) 如 \eqref{eq:xi-slice-carath-def}，
则“核版 (G)RH 在底 \(L\) 上成立”（逐通道临界线约束）等价于对所有通道 \(\rho\) 同时成立的视界正性/证书层级：
\[
\Re\mathscr{C}_{\Phi_{L,\rho}}(\omega)\ge 0\quad(\omega\in\mathbb{D})
\qquad\Longleftrightarrow\qquad
\mathbf{T}^{(\rho)}_N\succeq 0\quad(\forall N\ge 0),
\]
其中 \(\mathbf{T}^{(\rho)}_N\) 为 \(\mathscr{C}_{\Phi_{L,\rho}}\) 的 Toeplitz 主子矩阵。
因此，正文的“区间根/单位圆根/Schur--Cohn LMI”型谱证书与附录 \ref{subsubsec:app-horizon-weyl-carath-toeplitz} 的视界 Carath\'eodory--Toeplitz PSD 层级在同一视界正性框架下合并为统一证书语言。
\end{corollary}

\begin{remark}[6D 参数几何中的 PSD 可行域（接口性陈述）]\label{rem:xi-6d-parameter-spectrahedral-feasible}
当通道行列式的视界多极矩（或其有限模态粗粒化；参见命题 \ref{prop:xi-horizon-21mode-weyl-compression}）在某个 \(6\) 维参数立方体（例如由六个二值通道势的均值几何 \(\mathcal{R}_{\boldsymbol{\varphi}}\subset[0,1]^6\) 所组织，\S\ref{subsec:pom-potential-convex-geometry}）上实现为一组\emph{仿射}参数化时，
每一级 Toeplitz--PSD 约束 \(\mathbf{T}_N\succeq 0\) 即成为关于该 \(6\) 维参数的线性矩阵不等式（LMI），从而把“核版 (G)RH/拉马努金可行性”解释为参数立方体内的一个 spectrahedral 可行域族。
该几何接口为对称性筛选与优化（例如 Weyl-不变压缩下的低维参数化）提供了直接可审计的凸证书入口。
\end{remark}

\begin{remark}[Joukowski 压缩与有限维 toy-RH 字典的对齐]\label{rem:xi-joukowski-toy-rh-alignment}
命题 \ref{prop:xi-rh-interval-criterion} 的变量替换
\(
y=L^{s-1/2}
\)
与
\(
S=y+y^{-1}
\)
可视为 Joukowski 映射在 Dirichlet 坐标下的实现：它把单位圆 \(|y|=1\) 压缩到实区间 \([-2,2]\)，并把端点
\(
S=\pm 2
\)
对应到
\(
y=\pm 1
\)。
因此，“零点是否落在临界线 \(\Re(s)=\tfrac12\)”在 \(y\)-坐标中等价于“相应根是否落在单位圆”，而在 \(S\)-坐标中等价于“相应根是否落在 \([-2,2]\)”。
\par
该几何压缩与附录 \ref{app:real-input-40-finite-rh} 的有限维 \(\zeta\) 字典完全同型：对谱半径 \(\lambda=\rho(M)\) 的有限维核，
把 Dirichlet 变量写作 \(z=\lambda^{-s}\) 后，其临界带极点位置由归一化谱半径
\(
|\mu|/\sqrt{\lambda}
\)
控制，toy-RH（或 Ramanujan 半径界）等价于所有非平凡谱点满足 \(|\mu|\le \sqrt{\lambda}\)（命题 \ref{prop:finite-rh-triangle-dict} 与推论 \ref{cor:finite-rh-halfstrip-nopoles}）。
在“达界”情形 \(|\mu|=\sqrt{\lambda}\) 下，令
\(
y:=\mu/\sqrt{\lambda}
\)
则 \(|y|=1\) 且 \(S=y+y^{-1}\in[-2,2]\)，与本节的 \((y,S)\) 压缩坐标一致。
特别地，真实输入 \(40\) 状态核的唯一达界特征值为 \(-\sqrt{\lambda}\)（注 \ref{rem:finite-rh-40-ramanujan}），对应 \(y=-1\) 与 \(S=-2\)，从而“单频临界振荡”可被理解为该 Joukowski 压缩在端点的饱和（并见注 \ref{rem:finite-rh-40-hidden-minusA0} 对其结构来源的进一步分解）。
\end{remark}

\begin{remark}[半角提升的端点分解范式（与 Joukowski 压缩同型）]\label{rem:xi-half-angle-endpoint-paradigm}
除 \((y,S)=(L^{s-1/2},\,y+y^{-1})\) 的 Joukowski 压缩外，本文在 unitary twist 变量 \(u=e^{i\theta}\) 上还存在一个完全同型的“半角完成化”压缩：取半角变量 \(\mathfrak{r}=e^{i\theta/2}\) 使 \(u=\mathfrak{r}^2\)，并令 \(s=\mathfrak{r}+\mathfrak{r}^{-1}\in[-2,2]\)。
则幂半群 \(u\mapsto u^m\) 在完成化坐标中严格对应 \(s\mapsto C_m(s)\)（Chebyshev--Adams），并且端点 \(u=-1\)（\(s=0\)）处出现可审计的奇/偶通道喷射分裂：奇 \(m\) 通道控制线性主尺度，偶 \(m\) 通道只进入高阶喷射与有限部；同时 primitive 多项式的奇长度扇区在 \(u=-1\) 处严格消失，从而端点主奇性由偶长度扇区独占。
该端点分解的统一表述见附录 \ref{app:unit-circle-phase-gate} 的命题 \ref{prop:endpoint-decomposition-paradigm}。
\end{remark}
